% AI-written, AI cross-reviewed; no human sign-off; not a claim that RH is proved
% Atlas unit W04. Local manuscript v3, item 5.40, heading on printed PDF p. 49.
% UpperTheory.V.Z.Tf.SR.descWeak
% paper2/split_p013.tex:64–80. Verbatim excerpt; SHA-256 933a577eeedd803765f95d04a50f849bd9c43bedb855656df18973c0bc632641.
% Original TeX uses the macros and cross-reference keys of the source manuscript.

\begin{formalstatement}[Weak unitary descent in ranks two through four]\label{sub:W04}
Let $E/E^+$ be CM and unramified at finite places. Fix a prime $p$ and a
finite set $S$ of finite places of $E^+$ containing those above $p$,
all split in $E$. Let $2\le n\le4$ and let $\pi$ be a regular algebraic
cuspidal representation of $\mathrm{GL}_n(\mathbb A_E)$ with
$\pi^\vee\simeq\pi^c$ and Satake parameters at all places above the
complement of $S$. There exists an automorphic representation $\sigma$ of
the specified unitary group, unramified at every nonsplit finite place.
At all but finitely many split $w$ the component of $\sigma$, read
on $\mathrm{GL}_n(E_w)$, has the same Satake multiset as $\pi_w$.
\end{formalstatement}\subnode{W04}
Labesse's weak descent of a cohomological conjugate-dual cusp form
appears in Theorem~5.4. Its next assertion takes a member spherical at
all finite inert places as an additional hypothesis. W04 requires the
existence of a member with this property. In Lean, split-place rigidity
rules out the alternative constituent system in the trace comparison. Labesse also assumes $[E^+:\mathbb Q]\ge2$ and a unitary group quasi-split at inert finite places. As in W02, the finite-place unramifiedness assumption rules out $E^+=\mathbb Q$.
\href{https://webusers.imj-prg.fr/~michael.harris/CBCM.pdf}{Labesse, Theorem 5.4.}~\cite{r21-labesse}
